\SourcePage{19}{22}
\section{Operators}

Consider some physical quantity $f$ that characterizes the state of a quantum
system. Strictly speaking, in what follows we ought to speak not of a single
quantity but of an entire complete set of them at once. However, none of the
reasoning is thereby essentially altered, and for the sake of brevity and
simplicity we speak below of only a single physical quantity.

Those values which a given physical quantity may assume are termed its
\emph{eigenvalues} in quantum mechanics, and their aggregate is referred to as
the eigenvalue \emph{spectrum} of that quantity. In classical mechanics,
quantities range, in general, over a continuous series of values. In quantum
mechanics too there exist physical quantities (for example, coordinates) whose
eigenvalues fill a continuous range; in such cases one speaks of a
\emph{continuous spectrum} of eigenvalues. Alongside these quantities, however,
quantum mechanics also possesses quantities whose eigenvalues constitute a
discrete set; one then speaks of a \emph{discrete spectrum}.

For simplicity we shall first assume that the quantity $f$ under consideration
has a discrete spectrum; the case of a continuous spectrum is considered in
\secref{5}. We denote the eigenvalues of $f$ by $f_n$, where the index $n$ runs
through the values $0, 1, 2, 3, \ldots$ We denote the wave function of the
system in the state in which $f$ has the value $f_n$ by $\Psi_n$. The wave
functions $\Psi_n$ bear the name \emph{eigenfunctions} of the physical
quantity $f$. Every such function is taken to be normalized:
\begin{equation}
  \int |\Psi_n|^2\,dq = 1.
  \tag{3.1}
\end{equation}

If the system is in some arbitrary state with wave function $\Psi$, then a
measurement of $f$ performed on it will yield as a result one of the
eigenvalues $f_n$.%
\SourcePage{20}{23}
By the superposition principle, the wave
function $\Psi$ must be a linear combination of those eigenfunctions $\Psi_n$
that correspond to the values $f_n$ which can be found with non-zero probability
in a measurement performed on the system in the state under consideration.
In the general case of an arbitrary state, $\Psi$ may accordingly be written as the series
\begin{equation}
  \Psi = \sum_n a_n\Psi_n,
  \tag{3.2}
\end{equation}
where the summation is over all $n$, and $a_n$ are certain constant
coefficients.

Thus we arrive at the conclusion that every wave function can be, as one says,
expanded in the eigenfunctions of any physical quantity. A system of functions
in terms of which such an expansion can be carried out is spoken of as a
\emph{complete system of functions}.

The expansion (3.2) makes it possible to determine the probabilities of
finding (by means of measurements) one or another value $f_n$ of $f$ for a
system in the state with wave function $\Psi$. Indeed, according to what was
said in the preceding section, these probabilities must be determined by certain
expressions bilinear in $\Psi$ and $\Psi^*$ and must therefore be bilinear in
$a_n$ and $a_n^*$. Moreover, positivity of these expressions is obviously
required. As for the probability associated with a given $f_n$: it must
become unity for a system in the state $\Psi=\Psi_n$ and must vanish whenever
expansion (3.2) of $\Psi$ lacks a term containing the corresponding $\Psi_n$.
The only essentially positive quantity satisfying this condition is the squared
modulus of the coefficient $a_n$. Thus we arrive at the result that the squared
modulus $|a_n|^2$ of every expansion coefficient in (3.2) gives the
probability of the respective value $f_n$ of $f$ in the state with wave
function $\Psi$. The sum of the probabilities of all possible values $f_n$ must
equal unity; in other words, the relation
\begin{equation}
  \sum_n |a_n|^2 = 1.
  \tag{3.3}
\end{equation}
must hold.

If $\Psi$ were not normalized, then the relation (3.3) would also not hold. The
sum $\sum_n |a_n|^2$ would then have to be determined by some expression
bilinear in $\Psi$ and $\Psi^*$ and reducing to unity for normalized $\Psi$.
\SourcePage{21}{24}
The only such expression is the integral $\int\Psi\Psi^*\,dq$. Thus the
equality
\begin{equation}
  \sum_n a_n a_n^* = \int \Psi\Psi^*\,dq.
  \tag{3.4}
\end{equation}
must hold.

On the other hand, multiplying the expansion
$\Psi^*=\sum_n a_n^*\Psi_n^*$ of the complex conjugate function $\Psi^*$ by
$\Psi$ and integrating, we obtain
\[
  \int \Psi\Psi^*\,dq
  = \sum_n a_n^*\int \Psi_n^*\Psi\,dq.
\]
Comparing this with (3.4), we have
\[
  \sum_n a_n a_n^*
  = \sum_n a_n^*\int \Psi_n^*\Psi\,dq,
\]
from which we find the following formula determining the coefficients $a_n$ of
the expansion of $\Psi$ in the eigenfunctions $\Psi_n$:
\begin{equation}
  a_n = \int \Psi\Psi_n^*\,dq.
  \tag{3.5}
\end{equation}

Substituting (3.2) here, we obtain
\[
  a_n = \sum_m a_m\int \Psi_m\Psi_n^*\,dq,
\]
whence it is seen that the eigenfunctions must satisfy the conditions
\begin{equation}
  \int \Psi_m\Psi_n^*\,dq = \delta_{nm},
  \tag{3.6}
\end{equation}
where $\delta_{nm}=1$ for $n=m$ and $\delta_{nm}=0$ for $n\ne m$. The fact that
the integrals of the products $\Psi_m\Psi_n^*$ with $m\ne n$ vanish is spoken of
as the mutual \emph{orthogonality} of the functions $\Psi_n$. The collection of
eigenfunctions $\Psi_n$ thus constitutes a complete system of normalized and mutually
orthogonal (or, as one says for short,~\emph{orthonormal}) functions.

We introduce the concept of the \emph{mean value} $\bar f$ of a quantity $f$ in
a given state. Following the standard definition,
$\bar f$ is the sum of every eigenvalue $f_n$ of the quantity in question, weighted by the respective%
\SourcePage{22}{25}
probability $|a_n|^2$:
\begin{equation}
  \bar f = \sum_n f_n|a_n|^2.
  \tag{3.7}
\end{equation}

We write $\bar f$ in the form of an expression that involves not the expansion
coefficients of $\Psi$ but the function itself. Since (3.7) contains the
products $a_n a_n^*$, it is clear that the desired expression must be bilinear
in $\Psi$ and $\Psi^*$. We introduce a certain mathematical \emph{operator},
which we denote by $\hat f$\footnote{We shall everywhere denote operators by
letters with a hat.}, and define as follows. Let $(\hat f\Psi)$ denote the
result of the action of the operator $\hat f$ on the function $\Psi$. We define
$\hat f$ by requiring that the integral of $(\hat f\Psi)$ multiplied by the
conjugate $\Psi^*$ yield $\bar f$:
\begin{equation}
  \bar f = \int \Psi^*(\hat f\Psi)\,dq.
  \tag{3.8}
\end{equation}

It is easy to see that in the general case the operator $\hat f$ is some linear
integral operator.\footnote{A linear operator is one possessing the properties:
$\hat f(\Psi_1+\Psi_2)=\hat f\Psi_1+\hat f\Psi_2$,
$\hat f(a\Psi)=a\hat f\Psi$, where $\Psi_1$, $\Psi_2$ are arbitrary functions
and $a$ is an arbitrary constant.} Indeed, using the expression (3.5) for
$a_n$, the definition (3.7) of the mean value can be cast as
\[
  \bar f = \sum_n f_n a_n a_n^*
  = \int \Psi^*\left(\sum_n a_n f_n\Psi_n\right)dq.
\]
Comparing with (3.8), we see that the result of the action of the operator
$\hat f$ on the function $\Psi$ has the form
\begin{equation}
  (\hat f\Psi)=\sum_n a_n f_n\Psi_n.
  \tag{3.9}
\end{equation}
Substituting here the expression (3.5) for $a_n$, we find that $\hat f$ is an
integral operator of the form
\begin{equation}
  (\hat f\Psi)=\int K(q,q')\Psi(q')\,dq',
  \tag{3.10}
\end{equation}
where the function $K(q,q')$ (the so-called kernel of the operator) is
\begin{equation}
  K(q,q')=\sum_n f_n\Psi_n^*(q')\Psi_n(q).
  \tag{3.11}
\end{equation}
\SourcePage{23}{26}
Thus, to each physical quantity in quantum mechanics there is assigned a
definite linear operator.

From (3.9) it is seen that if $\Psi$ is one of the eigenfunctions $\Psi_n$ (so
that all $a_n$ except one are zero), then upon the action of the operator
$\hat f$ on it, this function is simply multiplied by the eigenvalue
$f_n$\footnote{Below, wherever it cannot lead to misunderstanding, we shall
omit the parentheses in the expression $(\hat f\Psi)$, the operator being
understood to act on the expression written after it.}:
\begin{equation}
  \hat f\Psi_n=f_n\Psi_n.
  \tag{3.12}
\end{equation}

The eigenfunctions of a physical quantity $f$ are therefore solutions of the
equation
\[
  \hat f\Psi=f\Psi,
\]
in this equation $f$ denotes a constant; the eigenvalues are those particular
values of $f$ that yield solutions obeying the necessary conditions. As will be seen later, the explicit form of operators for the various
physical quantities can be established on direct physical grounds, whereupon
this operator property enables one to obtain eigenfunctions and
eigenvalues by solving the equations $\hat f\Psi=f\Psi$.

For a real physical quantity, both the eigenvalues themselves and their mean
values in any state are real. This fact places a definite constraint on the
properties of the associated operators. Setting the expression (3.8) equal to
its own complex conjugate gives
\begin{equation}
  \int \Psi^*(\hat f\Psi)\,dq
  = \int \Psi(\hat f^*\Psi^*)\,dq,
  \tag{3.13}
\end{equation}
where $\hat f^*$ denotes the operator complex conjugate to
$\hat f$\footnote{By definition, if for the operator $\hat f$ we have
$\hat f\psi=\varphi$, then the complex-conjugate operator $\hat f^*$ is the
operator for which $\hat f^*\psi^*=\varphi^*$.}. For an arbitrary linear
operator such a relation does not, in general, hold, so that it represents a
certain restriction imposed on the possible form of the operators $\hat f$. For
an arbitrary operator $\hat f$ one can define, as one says, the
\emph{transposed} operator $\tilde{\hat f}$, defined%
\SourcePage{24}{27}
so that
\begin{equation}
  \int \Phi(\hat f\Psi)\,dq
  = \int \Psi(\tilde{\hat f}\Phi)\,dq,
  \tag{3.14}
\end{equation}
where $\Psi$, $\Phi$ are two different functions. Choosing as
$\Phi$ the conjugate function $\Psi^*$, comparison with (3.13)
reveals that
\begin{equation}
  \tilde{\hat f}=\hat f^*.
  \tag{3.15}
\end{equation}
An operator that satisfies this condition is called \emph{Hermitian}\footnote{For a
linear integral operator of the form (3.10), Hermiticity requires
that the kernel of the operator satisfy
$K(q,q')=K^*(q',q)$.}. Operators representing real physical quantities within
the mathematical apparatus of quantum mechanics must therefore be Hermitian.

One may also formally consider complex physical quantities,
i.e. quantities possessing complex eigenvalues. Suppose $f$ is one such
quantity; one may then form the complex-conjugate quantity $f^*$, whose eigenvalues are the
complex conjugates of the eigenvalues of $f$. The operator corresponding to the
quantity $f^*$ will be denoted by $\hat f^+$. It is called the \emph{adjoint}
operator to $\hat f$ and must, in general, be distinguished from the
complex-conjugate operator $\hat f^*$. Indeed, by definition of the operator
$\hat f^+$, the mean value of $f^*$ in some state $\Psi$ is
\[
  \overline{f^*}=\int \Psi^*\hat f^+\Psi\,dq.
\]
On the other hand, we have
\[
  {\left(\bar f\right)}^*
  = {\left[\int \Psi^*\hat f\Psi\,dq\right]}^*
  = \int \Psi\hat f^*\Psi^*\,dq
  = \int \Psi^*\tilde{\hat f}^{\,*}\Psi\,dq.
\]
Equating the two expressions, we find that
\begin{equation}
  \hat f^+=\tilde{\hat f}^{\,*},
  \tag{3.16}
\end{equation}
from which it is clear that $\hat f^+$ does not, in general, coincide with
$\hat f^*$. The condition (3.15) can now be written in the form
\begin{equation}
  \hat f=\hat f^+,
  \tag{3.17}
\end{equation}
\SourcePage{25}{28}
i.e. the operator of a real physical quantity coincides with its adjoint
(Hermitian operators are also called \emph{self-adjoint}).

Let us show how one can directly prove the mutual orthogonality of the
eigenfunctions of a Hermitian operator corresponding to different eigenvalues.
Let $f_n$, $f_m$ be two different eigenvalues of the real quantity $f$, and
$\Psi_n$, $\Psi_m$ the corresponding eigenfunctions:
\[
  \hat f\Psi_n=f_n\Psi_n,
  \qquad
  \hat f\Psi_m=f_m\Psi_m.
\]
Multiplying both sides of the first of these equalities by $\Psi_m^*$, and the
complex conjugate of the second equality by $\Psi_n$, and subtracting these
products termwise from each other, we obtain
\[
  \Psi_m^*\hat f\Psi_n-\Psi_n\hat f^*\Psi_m^*
  =(f_n-f_m)\Psi_n\Psi_m^*.
\]
Integrate both sides of this equality over $dq$. Since
$\hat f^*=\tilde{\hat f}$, by virtue of (3.14) the integral of the left-hand
side of the equality vanishes, so that we obtain
\[
  (f_n-f_m)\int\Psi_n\Psi_m^*\,dq=0,
\]
whence, in view of $f_n\ne f_m$, the required orthogonality property of the
functions $\Psi_n$ and $\Psi_m$ follows.

We have been speaking here all along of only a single physical quantity $f$,
whereas we ought to speak, as was noted at the beginning of the section, of a
complete system of simultaneously measurable physical quantities. Then we would
find that to each of these quantities $f$, $g$, $\ldots$ there corresponds its
own operator $\hat f$, $\hat g$, $\ldots$ The eigenfunctions $\Psi_n$
correspond to states in which all the quantities under consideration have
definite values, i.e. correspond to definite sets of eigenvalues
$f_n$, $g_n$, $\ldots$, and are joint solutions of the system of equations
\[
  \hat f\Psi=f\Psi,
  \qquad
  \hat g\Psi=g\Psi,
  \ldots
\]
